\chapter{Registre des paramètres}
\label{ann:registre}

Le registre complet (46 paramètres) est fourni dans \code{outputs/registre\_parametres.csv}, avec sa source, sa méthode,
sa période d'estimation et les mécanismes qui en dépendent. Catégories : E estimé sur données françaises, C calibré sur
données françaises, L transposé de la littérature, H hypothèse conventionnelle. La plage est celle utilisée dans les
analyses de sensibilité ($\pm$ 2 écarts-types pour les coefficients estimés).

\begin{scriptsize}
\begin{longtable}{@{}p{5.6cm}rcrlp{3.6cm}@{}}
\toprule
Paramètre & Valeur & Cat. & $t$ & Confiance & Plage \\
\midrule\endfirsthead
\toprule Paramètre & Valeur & Cat. & $t$ & Confiance & Plage \\ \midrule\endhead
\bottomrule\endfoot
Consommation : élasticité de court terme au RDB réel & 0{,}853 & E & 4.2 & élevée & 0.442 à 1.264 ($\pm$ 2 e.-t.) \\
Consommation : force de rappel & $-$0{,}218 & E & $-$1.0 & faible & -0.640 à 0.204 ($\pm$ 2 e.-t.) \\
Importations : élasticité de court terme à la demande & 1{,}46 & E & 15.9 & élevée & 1.281 à 1.648 ($\pm$ 2 e.-t.) \\
Importations : élasticité-prix de court terme & $-$0{,}114 & E & $-$1.1 & faible & -0.317 à 0.089 ($\pm$ 2 e.-t.) \\
Importations : prix relatif de long terme & $-$0{,}108 & E & $-$1.5 & faible & -0.254 à 0.039 ($\pm$ 2 e.-t.) \\
Importations : rappel & $-$0{,}661 & E & $-$4.4 & élevée & -0.960 à -0.362 ($\pm$ 2 e.-t.) \\
Exportations : élasticité à la demande mondiale & 0{,}711 & E & 7.5 & élevée & 0.522 à 0.900 ($\pm$ 2 e.-t.) \\
Exportations : élasticité-prix & $-$0{,}595 & E & $-$3.5 & élevée & -0.937 à -0.253 ($\pm$ 2 e.-t.) \\
Investissement des SNF : accélérateur & 1{,}95 & E & 5.8 & élevée & 1.279 à 2.612 ($\pm$ 2 e.-t.) \\
Investissement des SNF : rappel & $-$0{,}395 & E & $-$2.8 & moyenne & -0.674 à -0.117 ($\pm$ 2 e.-t.) \\
Investissement des SNF : coût d'usage du capital & $-$0{,}0728 & E & $-$3.0 & moyenne & -0.121 à -0.024 ($\pm$ 2 e.-t.) \\
Investissement par branche : accélérateur & 0{,}67 & E & 3.3 & élevée & 0.262 à 1.079 ($\pm$ 2 e.-t.) \\
Salaires : indexation de court terme & 0{,}394 & E & 1.8 & faible & -0.051 à 0.839 ($\pm$ 2 e.-t.) \\
Salaires : rappel vers la part salariale & $-$0{,}093 & E & $-$1.6 & faible & -0.207 à 0.020 ($\pm$ 2 e.-t.) \\
Prix de VA : transmission du coût salarial unitaire & 0{,}67 & E & 5.9 & élevée & 0.444 à 0.896 ($\pm$ 2 e.-t.) \\
Prix de VA : rappel & $-$0{,}117 & E & $-$2.0 & moyenne & -0.232 à -0.002 ($\pm$ 2 e.-t.) \\
Prix de VA : tension des capacités & 0{,}0015 & E & 1.7 & faible & -0.000 à 0.003 ($\pm$ 2 e.-t.) \\
Emploi : élasticité de court terme à la production & 0{,}259 & E & 6.0 & élevée & 0.172 à 0.346 ($\pm$ 2 e.-t.) \\
Emploi : rappel & $-$0{,}206 & E & $-$6.6 & élevée & -0.269 à -0.144 ($\pm$ 2 e.-t.) \\
Logement des ménages : élasticité au revenu & 1{,}48 & E & 4.2 & élevée & 0.783 à 2.185 ($\pm$ 2 e.-t.) \\
Logement des ménages : rappel & $-$0{,}282 & E & $-$1.4 & faible & -0.686 à 0.122 ($\pm$ 2 e.-t.) \\
Transmission des prix d'importation (année même) & 0{,}178 & E & 3.1 & élevée & 0.064 à 0.293 ($\pm$ 2 e.-t.) \\
Revenus de la propriété des ménages : élasticité au PIB nominal & 0{,}874 & E & 1.1 & faible & -0.648 à 2.395 ($\pm$ 2 e.-t.) \\
Revenu mixte et EBE des ménages : élasticité au PIB nominal & 0{,}89 & E & 4.7 & élevée & 0.509 à 1.271 ($\pm$ 2 e.-t.) \\
Courbe de Phillips (E.PHILLIPS) & $-$0{,}004 & C & — & faible & −0,008 à −0,002 \\
Élasticité de substitution capital-travail (SIGMA\_TRAVAIL) & 0{,}4 & L & — & faible & 0,2 à 0,7 \\
Fermeture d'offre (RHO\_DEFAUT) & 0{,}15 & H & — & faible & 0 à 0,4 \\
Loi d'Okun (OKUN) & 1{,}5 & L & — & moyenne & 1,2 à 2 \\
Délai de construction des capacités (DELAI\_CAPACITES) & 2 & H & — & faible & 1 à 4 ans \\
Profil des propensions à consommer par décile & 0,90 … 0,40 & L & — & faible & profil plat à accentué \\
Propension du dernier centile (PMC\_TOP1) & 0{,}25 & L & — & faible & 0,1 à 0,5 \\
Part distribuée des profits supplémentaires & 0{,}5 & L & — & faible & 0,3 à 0,7 \\
Part des dividendes aux non-résidents & 0{,}3 & C & — & moyenne & 0,20 à 0,50 \\
Part des dividendes résidents reçue par les ménages & 0{,}35 & C & — & faible & 0,2 à 0,5 \\
Élasticité du revenu imposable (M13) & 0{,}2 & L & — & moyenne & 0,1 à 0,4 \\
Élasticité des revenus du capital (M13) & 0{,}5 & L & — & moyenne & 0,2 à 1 \\
Perte d'assiette de l'ISF (M13) & 0{,}25 & L & — & faible & 0,1 à 0,5 \\
Élasticité de la production au capital public & 0{,}08 & L & — & moyenne & 0,04 à 0,12 \\
Coût public d'un point de SMIC & 660 & C & — & moyenne & 500 à 800 M€ \\
Diffusion du SMIC au salaire moyen & 0{,}11 & H & — & faible & 0,05 à 0,2 \\
Répercussion de la TVA sur les prix & 0{,}8 & L & — & moyenne & 0,6 à 1 \\
Répercussion des impôts de production & 0{,}5 & H & — & faible & 0,2 à 0,8 \\
Compensation horaire des 32 heures & 0{,}7 & H & — & faible & 0,5 à 0,9 \\
Coefficient de pension marginale & 0{,}55 & C & — & moyenne & 0,45 à 0,7 \\
Taux d'épargne cible de la référence & 17,9 → 15,5 \% & H & — & faible & 14 à 18 \% en 2031 \\
Rétroactions zone euro (THETA\_PRIME, etc.) & M11 & C & — & faible & 0 à 2 × central \\

\end{longtable}
\end{scriptsize}

\chapter{Glossaire}
\label{ann:glossaire}

\begin{description}[style=nextline,leftmargin=1em,font=\normalfont\bfseries\color{bleu}]
\small
\item[Ajustement stock-flux] Écart entre la variation de la dette publique et le déficit, dû aux opérations financières
(prêts, cessions, trésorerie) et aux effets de valorisation. Nul en projection dans le modèle.
\item[Calage, calibration] Fixation de la valeur d'un paramètre ou d'un terme pour que le modèle reproduise des
grandeurs observées, sans estimation statistique du paramètre lui-même.
\item[Chômage structurel] Taux de chômage compatible avec une inflation stable ; dans le modèle, 7,5~\% (hypothèse), cible
de la fermeture d'offre.
\item[Coefficient technique] Quantité d'un produit nécessaire pour produire une unité d'un autre produit (tableau
entrées-sorties).
\item[Contenu en importations] Part de la valeur d'une demande finale qui correspond, directement ou indirectement, à des
importations.
\item[Cotisations sociales] Prélèvements sur les salaires (part salariale et part employeur) qui financent la protection
sociale ; les allègements généraux en réduisent le montant sur les bas salaires.
\item[Coût d'usage du capital] Coût annuel de détention d'une unité de capital : taux d'intérêt plus dépréciation, corrigé de
l'impôt sur les sociétés, $(r+\delta)/(1-\tau)$.
\item[Coût salarial unitaire (CSU)] Coût du travail par unité produite : salaire par tête $\times$ emploi / production.
\item[Décile] Chacun des dix groupes de même effectif obtenus en classant la population selon son niveau de vie ; D1 et D9
désignent les seuils du premier et du neuvième décile.
\item[Écart-type] Mesure de la précision d'un coefficient estimé ; le rapport $t$ (coefficient / écart-type) supérieur à 2
indique en général un coefficient significativement différent de zéro.
\item[Effet boule de neige] Hausse spontanée du ratio de dette quand le taux d'intérêt apparent dépasse la croissance du
PIB nominal : $\Delta d \approx (i-g)\,d - s$, avec $s$ le solde primaire.
\item[Élasticité] Variation en pourcentage d'une grandeur provoquée par une variation de 1~\% d'une autre.
\item[Élasticité de substitution capital-travail ($\sigma$)] Variation en \% du rapport capital / travail quand leur prix
relatif varie de 1~\%.
\item[Équation à correction d'erreur (MCE)] Équation qui combine une réaction de court terme et un rappel progressif vers une
relation d'équilibre de long terme (encadré 2.1).
\item[Fermeture d'offre] Mécanisme du modèle qui ramène progressivement le chômage vers son niveau structurel par un
supplément de demande ; résume la politique monétaire, le crédit et la compétitivité de moyen terme. Hypothèse $\rho = 0{,}15$.
\item[Indice de Gini] Mesure de l'inégalité d'une distribution, de 0 (égalité parfaite) à 1 (une personne détient tout).
\item[Leontief (modèle de)] Modèle entrées-sorties qui calcule la production nécessaire pour satisfaire une demande
finale, compte tenu des achats entre branches : $x = (I-A)^{-1} f$.
\item[Loi d'Okun] Relation entre l'écart de production et le chômage : environ 1,5 point de PIB par point de chômage.
\item[Méta-modèle] Approximation d'un modèle complexe par une fonction simple (ici linéaire) estimée sur ses simulations,
utilisée pour l'optimisation.
\item[Microsimulation] Méthode qui applique la législation à un échantillon de ménages réels pour calculer impôts et
prestations ménage par ménage.
\item[Multiplicateur] Variation du PIB provoquée par une variation d'un euro de dépense publique ou de recette.
\item[Niveau de vie] Revenu disponible du ménage divisé par le nombre d'unités de consommation (1 pour le premier adulte,
0,5 pour les autres personnes de 14 ans ou plus, 0,3 pour les enfants de moins de 14 ans).
\item[Pauvreté monétaire] Situation des personnes dont le niveau de vie est inférieur à 60~\% du niveau de vie médian.
\item[PIB nominal, PIB en volume] Le PIB nominal (en valeur) est mesuré aux prix courants ; le PIB en volume élimine
l'effet des prix.
\item[Prime de risque] Supplément de taux d'intérêt exigé par les prêteurs quand la dette augmente ; hypothèse de 3 points
de base par point de dette.
\item[Programme linéaire] Problème d'optimisation dont l'objectif et les contraintes sont linéaires ; résolu ici par la
méthode HiGHS.
\item[Propension marginale à consommer] Part d'un euro de revenu supplémentaire consacrée à la consommation.
\item[Revenu disponible brut (RDB)] Revenus des ménages après impôts et cotisations et y compris prestations ; le
\emph{pouvoir d'achat du RDB} est le RDB divisé par le prix de la consommation.
\item[Revenu national brut (RNB)] PIB plus revenus primaires nets reçus du reste du monde (dividendes, intérêts, salaires
transfrontaliers).
\item[Seuil fixe (pauvreté au)] Taux de pauvreté calculé après réforme avec le seuil de pauvreté d'avant réforme, pour isoler
l'effet de la mesure.
\item[Solde public, solde primaire] Différence entre recettes et dépenses publiques ; le solde primaire exclut la charge
d'intérêts.
\item[Taux marginal effectif] Part d'une hausse de revenu brut reprise par les impôts, les cotisations et la baisse des
prestations sous condition de ressources.
\item[Taux d'utilisation des capacités (TUC)] Rapport entre la production et la production maximale permise par les
équipements.
\item[Termes de rappel ancrés] Dans la référence, l'écart à l'équilibre constaté fin 2025 est traité comme durable pour
éviter un rattrapage artificiel après 2025.
\item[Variante] Simulation d'une mesure comparée à la référence ; les résultats sont des écarts.
\end{description}

\chapter{Reproductibilité}
\label{ann:repro}

\paragraph{Code et données.} Dépôt du projet (dossier \code{rapport/}) :
\begin{itemize}
\item \code{rapport/simulations.py} : les 60 expériences des chapitres~\ref{ch:experiences} et~\ref{ch:complementaires}
(environ 90 minutes sur deux processeurs) ; résultats \code{rapport/resultats/experiences\_*.json} ;
\item \code{rapport/seuils\_salaires.py} : cas types OpenFisca ; \code{rapport/resultats/seuils\_salaires.json} ;
\item \code{rapport/strategie.py} : méta-modèle et programmes linéaires ; \code{rapport/resultats/strategie.json} ;
\item \code{rapport/verification.py} : vérification des scénarios avec le modèle complet, mise en œuvre partielle et
tests de robustesse ; \code{rapport/resultats/verification\_\{A,B,C,D,R\}.json} ;
\item \code{rapport/figures.py} : toutes les figures (\code{rapport/figures/}) et tous les tableaux générés
(\code{rapport/tableaux/}) ;
\item \code{rapport/tex/} : sources LaTeX de ce rapport (compilation : \code{xelatex rapport.tex}, deux passes) ;
\item \code{rapport/livrables/} : synthèse des mesures et de leurs impacts (\code{mesures\_impacts.csv},
\code{.xlsx}), paramètres des scénarios (\code{scenarios\_parametres.json}), résultats détaillés (\code{resultats\_detailles.csv}).
\end{itemize}

\paragraph{Version.} Modèle 0.6 (dépôt git du projet). La version 0.5 reste reproductible (commit \code{6a84ecd}).
Les données sont téléchargées par script depuis les sources ouvertes et mises en cache ; le registre des téléchargements
(\code{outputs/registre\_telechargements.json}) conserve les adresses et les dates.

\paragraph{Données utilisées.}
\begin{small}
\begin{tabularx}{\linewidth}{@{}p{4.2cm}L@{}}
\toprule
Source & Séries \\
\midrule
Insee & Comptes nationaux annuels (base 2020, mai 2026), comptes des administrations publiques, enquête Emploi (chômage BIT),
projections de population 2026, recensement 2022 (fichier détail Individus), indices de prix, enquête de conjoncture
(taux d'utilisation des capacités) \\
Eurostat & Tableaux entrées-sorties (\code{naio\_10\_cp1700}, \code{naio\_10\_pyp1700}, \code{naio\_10\_cp15}), comptes
de branches (\code{nama\_10\_a64}, \code{nama\_10\_a64\_e}, \code{nama\_10\_a64\_p5}, \code{nama\_10\_nfa\_st}), comptes de
secteurs (\code{nasa\_10\_nf\_tr}), finances publiques (\code{gov\_10a\_main}, \code{gov\_10dd\_edpt1}), taux longs
(\code{irt\_lt\_mcby\_a}), activité (\code{lfsa\_argan}), pensions (\code{spr\_exp\_pens}, \code{spr\_pns\_ben}) \\
DGFiP & IRCOM (revenus 2024), statistiques nationales par case de déclaration 2042 \\
Drees, Cnaf & Pensions moyennes, montants et effectifs des prestations \\
Banque de France & Détention des actions par les non-résidents, balance des paiements \\
FMI & Perspectives de l'économie mondiale (octobre 2026), API DataMapper \\
COR & Rapport annuel (juin 2026) : hypothèses de productivité et de chômage \\
OpenFisca-France & Législation socio-fiscale (version 177) \\
\bottomrule
\end{tabularx}
\end{small}
